% ============================================================
%  Tiểu luận Nghiên cứu Khoa học
%  Đề tài: Fine-tuning Qwen2.5-1.5B với QLoRA
%          cho các Tác vụ Suy luận Pháp lý Tiếng Việt
%  Sinh viên: Lương Vân Khoa, Doãn Sơn Hoàng, Bàn Khánh Duy
% ============================================================

\documentclass[fontsize=13pt,a4paper,oneside,DIV=calc,chapterprefix=true]{scrbook}

\usepackage[utf8]{vietnam}
\usepackage{husthesis}
\usepackage{enumitem}
\usepackage{tipa}
\usepackage{amssymb}
\usepackage{graphicx}
\usepackage{subcaption}
\usepackage{booktabs}
\usepackage{mathptmx}
\usepackage{amsmath}
\usepackage{amsfonts}
\usepackage{array}
\newcolumntype{P}[1]{>{\centering\arraybackslash}p{#1}}
\newcolumntype{M}[1]{>{\centering\arraybackslash}m{#1}}
\usepackage{fancyhdr}
\usepackage{multirow}
\usepackage{float}
\usepackage{longtable}
\usepackage{caption}
\usepackage{listings}
\usepackage{xcolor}
\usepackage{hyperref}

\hypersetup{
    colorlinks=true,
    linkcolor=blue,
    filecolor=magenta,
    urlcolor=cyan,
    citecolor=blue,
}

\lstset{
  basicstyle=\ttfamily\footnotesize,
  breaklines=true,
  frame=single,
  backgroundcolor=\color{gray!8},
  keywordstyle=\color{blue}\bfseries,
  commentstyle=\color{green!50!black},
  stringstyle=\color{red!60!black},
  numbers=left,
  numberstyle=\tiny\color{gray},
  numbersep=5pt,
  showstringspaces=false,
}

% ---- Thông tin luận văn ----
\crname{Tiểu luận Nghiên cứu Khoa học}
\ctname{FINE-TUNING QWEN2.5-1.5B VỚI QLORA\\
        CHO CÁC TÁC VỤ SUY LUẬN PHÁP LÝ TIẾNG VIỆT}
\cstuname{
  Lương Vân Khoa \\[4pt]
  Doãn Sơn Hoàng \\[4pt]
  Bàn Khánh Duy \\
}
\csMajor{Khoa học máy tính và thông tin}
\csProgram{chuẩn}
\csSupervise{TS. ---}
\csdeptname{KHOA TOÁN -- CƠ -- TIN HỌC}
\cttime{2026}

\thesislayout

\begin{document}

\coverpage
\secondcoverpage

\frontmatter

% ---- Lời cam đoan ----
\begin{declaration}
Nhóm sinh viên thực hiện tiểu luận này cam đoan rằng mọi nội dung trình bày trong tài liệu
là kết quả của quá trình nghiên cứu độc lập, dưới sự hướng dẫn của giảng viên. Các số liệu
thực nghiệm được ghi lại trung thực từ quá trình chạy thực tế trên nền tảng Google Colab;
không có số liệu nào được bịa đặt hoặc sao chép từ công trình khác mà không có trích dẫn
rõ ràng. Mọi tài liệu tham khảo đều được liệt kê đầy đủ ở cuối báo cáo.

\vspace{1cm}
\hfill\textit{Hà Nội, tháng 4 năm 2026}\\[0.5cm]
\hfill Nhóm sinh viên thực hiện
\end{declaration}

% ---- Lời cảm ơn ----
\begin{acknowledgments}
Trước tiên, nhóm xin gửi lời cảm ơn đến giảng viên hướng dẫn đã định hướng, góp ý và
kiên nhẫn theo dõi quá trình thực hiện đề tài từ những bước đầu tiên.

Ban tổ chức cuộc thi VLSP 2025 đã công bố bộ dữ liệu \textit{LegalSLM Public-Test} một cách
miễn phí, tạo điều kiện thuận lợi cho sinh viên và các nhóm nghiên cứu có ít tài nguyên tiếp
cận được dữ liệu pháp lý chất lượng cao. Cộng đồng HuggingFace, đặc biệt là các dự án
Transformers, PEFT và TRL, đã cung cấp những công cụ mã nguồn mở giúp rút ngắn đáng kể thời
gian triển khai.

Cuối cùng, nhóm cảm ơn Alibaba Cloud đã công bố mã nguồn và trọng số của dòng mô hình
Qwen2.5 theo giấy phép mở, giúp nghiên cứu có thể được thực hiện mà không cần chi phí
thương mại.
\end{acknowledgments}

% ---- Tóm tắt ----
\begin{abstract}
Bài báo này trình bày một nghiên cứu thực nghiệm về việc tinh chỉnh mô hình ngôn ngữ nhỏ
\textbf{Qwen2.5-1.5B-Instruct} cho ba tác vụ suy luận pháp lý tiếng Việt trong khuôn khổ
cuộc thi VLSP~2025 LegalSLM. Kỹ thuật được áp dụng là \textbf{QLoRA} --- kết hợp lượng
tử hóa 4-bit (NF4) với LoRA rank-16 --- cho phép huấn luyện hiệu quả trên GPU T4 (15~GB)
với chưa đến 2\% tham số được cập nhật.

Điểm đáng chú ý của phương pháp là \textbf{chiến lược huấn luyện đa tác vụ}: thay vì train
riêng lẻ từng tác vụ (như MinLegal~\cite{minlegal2025}), toàn bộ 120 mẫu từ ba tác vụ ---
Suy diễn ngôn ngữ tự nhiên (NLI), Trả lời câu hỏi trắc nghiệm (MCQ) và Suy luận tam đoạn
luận (Syllogism) --- được gộp thành một tập huấn luyện duy nhất. Cách tiếp cận này không chỉ
đơn giản hóa quy trình mà còn cho phép mô hình học cách phân biệt và xử lý nhiều loại câu
hỏi pháp lý qua cùng một LoRA adapter.

Kết quả thực nghiệm cho thấy cải thiện đáng kể so với baseline zero-shot trên cả ba tác vụ:
độ chính xác NLI tăng từ 66,67\% lên \textbf{93,33\%} (+40,0\%), MCQ từ 73,33\% lên
\textbf{86,67\%} (+18,2\%), ROUGE-L Syllogism từ 0,4230 lên \textbf{0,5066} (+19,8\%),
với toàn bộ quá trình huấn luyện hoàn thành trong \textbf{19,3 phút} trên GPU T4 miễn phí
của Google Colab --- chi phí tính toán cực kỳ thấp so với các phương pháp tham chiếu.
\end{abstract}

% ---- Danh sách từ viết tắt ----
\nomenclature{LLM}{Large Language Model --- Mô hình ngôn ngữ lớn}
\nomenclature{SLM}{Small Language Model --- Mô hình ngôn ngữ nhỏ}
\nomenclature{NLI}{Natural Language Inference --- Suy diễn ngôn ngữ tự nhiên}
\nomenclature{MCQ}{Multiple Choice Question --- Câu hỏi trắc nghiệm}
\nomenclature{QLoRA}{Quantized Low-Rank Adaptation}
\nomenclature{LoRA}{Low-Rank Adaptation}
\nomenclature{PEFT}{Parameter-Efficient Fine-Tuning --- Tinh chỉnh hiệu quả tham số}
\nomenclature{MTL}{Multi-Task Learning --- Huấn luyện đa tác vụ}
\nomenclature{SFT}{Supervised Fine-Tuning --- Tinh chỉnh có giám sát}
\nomenclature{ROUGE}{Recall-Oriented Understudy for Gisting Evaluation}
\nomenclature{NF4}{NormalFloat 4-bit --- Định dạng 4-bit phân bố chuẩn}
\nomenclature{VRAM}{Video RAM --- Bộ nhớ GPU}
\nomenclature{VLSP}{Vietnamese Language and Speech Processing}
\printnomenclature

\tableofcontents
\listoftables
\listoffigures

% ============================================================
\mainmatter

\fancyhead{}
\renewcommand{\footrulewidth}{0.4pt}
\pagestyle{fancy}
\fancyhead[RO]{\small\thepage}
\fancyhead[LO]{\small\itshape Fine-tuning Qwen2.5-1.5B với QLoRA}
\fancyfoot[C]{}

% ============================================================
\chapter*{Giới thiệu}
\label{chap:intro}
\addcontentsline{toc}{chapter}{\nameref{chap:intro}}

\section*{Đặt vấn đề}

Hệ thống pháp luật Việt Nam ngày càng trở nên phức tạp với hàng chục nghìn văn bản quy phạm
pháp luật được ban hành và sửa đổi mỗi năm. Người dân, doanh nghiệp và cả các chuyên gia pháp
lý thường gặp khó khăn trong việc tra cứu, đối chiếu và suy luận từ các điều khoản luật.
Đây chính là bài toán mà các hệ thống Trí tuệ nhân tạo (AI) có thể hỗ trợ đáng kể.

Các mô hình ngôn ngữ lớn (LLM) như GPT-4 hay Gemini đã thể hiện khả năng ấn tượng trong hiểu
văn bản pháp lý, nhưng việc sử dụng chúng thường kéo theo hai trở ngại: chi phí API cao khi
triển khai quy mô lớn và sự thiếu chuyên biệt hóa cho ngữ cảnh pháp lý tiếng Việt. Xu hướng
nghiên cứu gần đây chuyển sang \textit{Small Language Models} (SLM) --- các mô hình dưới
4~tỷ tham số --- có thể chạy trên phần cứng thông thường và được tinh chỉnh chuyên sâu cho
một lĩnh vực cụ thể~\cite{bosch2025, minlegal2025}.

Trong khuôn khổ cuộc thi \textbf{VLSP 2025 LegalSLM}~\cite{vlsp2025}, các nhóm nghiên cứu
được yêu cầu phát triển mô hình ngôn ngữ nhỏ ($\leq$4B tham số) có khả năng giải quyết ba
tác vụ pháp lý tiếng Việt: (1) phân loại tính hữu ích của văn bản pháp lý (NLI), (2) trả lời
câu hỏi trắc nghiệm pháp luật (MCQ), và (3) lập luận pháp lý theo chuỗi tam đoạn luận
(Syllogism). Đây là bài toán khó do đặc thù ngôn ngữ pháp lý tiếng Việt vừa dài, vừa có
nhiều khái niệm chuyên ngành và các phụ thuộc logic phức tạp.

\section*{Mục tiêu và Phạm vi}

Nghiên cứu này đặt ra hai mục tiêu chính:
\begin{enumerate}
  \item \textbf{Kỹ thuật}: Chứng minh rằng kỹ thuật QLoRA kết hợp với chiến lược huấn luyện
    đa tác vụ có thể cải thiện đáng kể hiệu năng của mô hình ngôn ngữ nhỏ (1,5B tham số) trên
    ba tác vụ pháp lý tiếng Việt, chỉ với 120 mẫu huấn luyện và trong thời gian dưới 30 phút
    trên GPU T4 miễn phí.
  \item \textbf{Thực tiễn}: Cung cấp một pipeline hoàn chỉnh, có thể tái tạo trên Google Colab,
    làm cơ sở tham khảo cho các nghiên cứu tiếp theo về mô hình ngôn ngữ pháp lý tiếng Việt.
\end{enumerate}

Phạm vi của nghiên cứu được giới hạn ở việc đánh giá trên bộ dữ liệu VLSP~2025 LegalSLM
Public-Test, không mở rộng sang các bộ dữ liệu pháp lý khác. Nhóm không thực hiện pre-training
miền (domain pre-training) do hạn chế tài nguyên.

\section*{Những đóng góp chính}

\begin{enumerate}
  \item Trình bày phương pháp tinh chỉnh QLoRA đa tác vụ cho mô hình Qwen2.5-1.5B trên ba tác
    vụ NLI, MCQ và Syllogism bằng tiếng Việt, đạt kết quả cải thiện từ 18\% đến 40\% so với
    zero-shot baseline.
  \item Phân tích tác động của instruction tuning đa tác vụ, cho thấy mô hình học được cách
    phân biệt ba loại câu hỏi pháp lý mà không cần nhiều dữ liệu.
  \item Cung cấp notebook Google Colab có thể chạy hoàn toàn trên tài nguyên miễn phí, phù hợp
    với sinh viên và nhà nghiên cứu không có GPU chuyên dụng.
\end{enumerate}

\section*{Cấu trúc báo cáo}

Phần còn lại của báo cáo được tổ chức như sau: Chương~\ref{chap:background} trình bày cơ sở
lý thuyết và các công trình liên quan. Chương~\ref{chap:data} mô tả bộ dữ liệu và quy trình
tiền xử lý. Chương~\ref{chap:method} trình bày phương pháp đề xuất. Chương~\ref{chap:results}
báo cáo kết quả thực nghiệm và phân tích. Phần Kết luận tóm tắt và đề xuất hướng phát triển.

% ============================================================
\chapter{Cơ sở Lý thuyết và Công trình Liên quan}
\label{chap:background}

\section{Mô hình Ngôn ngữ Nhỏ (Small Language Model)}

Kể từ khi GPT-3 (175B tham số) được công bố vào năm 2020, cuộc đua xây dựng các mô hình ngày
càng lớn hơn đã kéo dài gần 3 năm. Tuy nhiên từ năm 2023, một xu hướng đối lập dần xuất hiện:
nhiều nhóm nghiên cứu tập trung vào các mô hình nhỏ hơn nhiều nhưng được huấn luyện kỹ lưỡng
hơn, thường được gọi là \textit{Small Language Model} (SLM).

SLM thường được định nghĩa là các mô hình dưới 7 tỷ tham số, đôi khi chặt chẽ hơn là dưới
4 tỷ. Các đại diện tiêu biểu bao gồm Phi-3.5-mini (3.8B, Microsoft), Gemma-2-2B (Google),
LLaMA-3.2-3B (Meta) và Qwen2.5 (0.5B đến 72B, Alibaba). So với các mô hình lớn, SLM có những
ưu điểm thực tiễn rõ ràng: chi phí suy diễn thấp, có thể chạy offline trên laptop phổ thông
hoặc điện thoại thông minh, và thời gian fine-tuning ngắn hơn nhiều. Đổi lại, năng lực lý luận
phức tạp và hiểu biết đa ngôn ngữ thường kém hơn so với các mô hình lớn.

\section{LoRA và Tinh chỉnh Hiệu quả Tham số}

\subsection{Bài toán full fine-tuning và vấn đề tài nguyên}

Tinh chỉnh toàn bộ tham số mô hình (full fine-tuning) yêu cầu gradient và optimizer states
cho \textit{tất cả} tham số. Với mô hình 1.5B tham số dùng AdamW, bộ nhớ cần thiết ước tính:
\begin{itemize}
  \item Trọng số (float16): $1.5 \times 10^9 \times 2\,\text{byte} \approx 3\,\text{GB}$
  \item Gradient: thêm $\approx 3\,\text{GB}$
  \item Adam states (m, v): thêm $\approx 12\,\text{GB}$
  \item Tổng cộng: $\approx 18\,\text{GB}$
\end{itemize}

Con số này vượt quá VRAM của GPU T4 (15~GB). Đây là lý do tại sao các kỹ thuật PEFT ra đời.

\subsection{LoRA: Low-Rank Adaptation}

LoRA~\cite{lora} xuất phát từ nhận xét rằng các cập nhật trọng số trong quá trình fine-tuning
thường có \textit{hạng nội tại thấp} (intrinsic low rank). Thay vì huấn luyện trực tiếp ma
trận trọng số $W \in \mathbb{R}^{d \times k}$, LoRA biểu diễn cập nhật dưới dạng tích hai
ma trận hạng thấp:
\begin{equation}
  W' = W_0 + \Delta W = W_0 + BA
  \label{eq:lora}
\end{equation}
trong đó $B \in \mathbb{R}^{d \times r}$, $A \in \mathbb{R}^{r \times k}$, và hạng $r \ll
\min(d, k)$. Trong quá trình huấn luyện, $W_0$ được đóng băng (frozen) hoàn toàn, chỉ $A$
và $B$ được cập nhật. Khi suy diễn, có thể hợp nhất $W_0 + BA$ thành một ma trận duy nhất
nên không tăng độ trễ suy diễn.

Số tham số huấn luyện giảm từ $d \times k$ xuống $r \times (d + k)$. Với $r = 16$ và
$d = k = 2048$ (một lớp attention thông thường), tỷ lệ nén là:
\begin{equation}
  \frac{r(d+k)}{dk} = \frac{16 \times 4096}{2048 \times 2048} \approx 1.56\%
\end{equation}

Ký hiệu $\alpha$ trong LoRA là hệ số tỷ lệ: đầu ra của adapter được nhân với $\frac{\alpha}{r}$
trước khi cộng vào $W_0 x$. Với $\alpha = 2r$, đây tương đương với learning rate effective gấp
đôi cho adapter.

\subsection{QLoRA: Quantized LoRA}

QLoRA~\cite{qlora} cải tiến LoRA theo hai hướng chính:

\textbf{(1) Lượng tử hóa NF4 (NormalFloat-4bit):} Mô hình nền được lưu trữ ở 4-bit thay vì
16-bit. Không dùng định dạng 4-bit thông thường (INT4) mà dùng \textit{NormalFloat-4 (NF4)} ---
một định dạng được thiết kế tối ưu cho phân bố Gaussian của trọng số neural network. NF4 chia
đều xác suất trên 16 điểm lượng tử hóa, giảm sai số lượng tử hóa so với INT4.

\textbf{(2) Double quantization:} Các hằng số lượng tử hóa (quantization constants) bản thân
cũng được lượng tử hóa lần thứ hai từ float32 xuống float8, tiết kiệm thêm khoảng 0.37~bit/tham
số trung bình.

Kết hợp hai kỹ thuật này, bộ nhớ cho mô hình 1.5B giảm từ $\approx 3$~GB xuống chỉ
$\approx 1.1$~GB. Trong khi đó, các ma trận LoRA vẫn được lưu và tính toán ở float16/bfloat16
để đảm bảo độ chính xác gradient. Tóm lại, QLoRA cho phép fine-tune mô hình 7B trên GPU 24~GB
hoặc mô hình 1.5B trên GPU 6~GB --- điều gần như không thể với full fine-tuning.

\section{Huấn luyện Đa Tác vụ}

Huấn luyện đa tác vụ (Multi-Task Learning, MTL) là phương pháp huấn luyện một mô hình đồng
thời trên nhiều tác vụ, kỳ vọng rằng các tác vụ sẽ chia sẻ thông tin hữu ích cho nhau~\cite{mtl_survey}.
Trong bối cảnh instruction tuning cho LLM, MTL thường được thực hiện bằng cách gộp dữ liệu
của các tác vụ vào một tập huấn luyện chung, với định dạng prompt khác nhau cho từng loại tác
vụ~\cite{flan2022}. Điều này tạo ra một mô hình ``generalist'' có thể phân biệt và xử lý
nhiều loại nhiệm vụ qua hướng dẫn (instruction).

Một điểm quan trọng cần lưu ý: MTL không phải lúc nào cũng vượt trội so với train riêng lẻ.
Khi các tác vụ quá khác biệt về cấu trúc đầu ra (ví dụ kết hợp phân loại với sinh văn bản
tự do), mô hình có thể gặp \textit{task interference} --- tác vụ này can thiệp vào việc học
tác vụ kia. Tuy nhiên, khi dữ liệu ít (như trong nghiên cứu này với chỉ 40 mẫu mỗi tác vụ),
MTL giúp tăng tổng số mẫu và giảm nguy cơ overfitting.

\section{Thước đo ROUGE-L}

ROUGE (Recall-Oriented Understudy for Gisting Evaluation) là bộ thước đo tiêu chuẩn đánh giá
chất lượng văn bản sinh ra so với tham chiếu, được dùng rộng rãi trong tóm tắt văn bản và
dịch máy. ROUGE-L cụ thể dựa trên \textit{Longest Common Subsequence} (LCS):
\begin{align}
  P_{lcs} &= \frac{LCS(X,Y)}{|Y|}, \quad
  R_{lcs} = \frac{LCS(X,Y)}{|X|} \\
  F_{lcs} &= \frac{(1+\beta^2) P_{lcs} R_{lcs}}{R_{lcs} + \beta^2 P_{lcs}}
\end{align}
Với $\beta = 1$ (F1 cân bằng), $X$ là câu tham chiếu, $Y$ là câu dự đoán. ROUGE-L đo độ
trùng lặp chuỗi không nhất thiết liên tục, phù hợp với văn bản pháp lý dài có nhiều biến thể
cách diễn đạt.

Hạn chế của ROUGE-L là chỉ đo trùng lặp từ ngữ bề mặt, không đánh giá được chất lượng lập
luận pháp lý. Một câu trả lời đúng về logic nhưng dùng từ khác nhau sẽ bị đánh giá thấp,
trong khi một câu trả lời sai về nội dung nhưng dùng nhiều từ giống tham chiếu lại bị đánh
giá cao. Đây là hạn chế đã biết của thước đo này với tác vụ sinh văn bản pháp lý.

\section{Các Công trình Liên quan}

\subsection{Bộ dữ liệu pháp lý tiếng Việt}

Nghiên cứu về NLP pháp lý tiếng Việt còn khá hạn chế so với tiếng Anh. Bộ dữ liệu VLQA~\cite{vlqa2025}
là tài nguyên lớn nhất hiện tại với hơn 3.000 cặp câu hỏi-đáp án được chú thích thủ công bởi
chuyên gia pháp lý, bao phủ 59.636 điều khoản thuộc 27 lĩnh vực. Cuộc thi VLSP~2025 LegalSLM
cung cấp thêm bộ Public-Test gồm 150 mẫu thuộc ba tác vụ, được lấy từ các câu hỏi pháp lý
thực tế đang lưu hành trên internet.

\subsection{Bosch@AI: Phương pháp đứng đầu VLSP 2025}

Nhóm Bosch@AI~\cite{bosch2025} đạt kết quả cao nhất tại VLSP~2025 LegalSLM với pipeline gồm
ba giai đoạn: (1) domain pre-training trên corpus pháp lý tiếng Việt thu thập từ nhiều nguồn,
(2) tổng hợp dữ liệu huấn luyện theo các ``aspect'' cụ thể của từng tác vụ bằng cách dùng LLM
lớn hơn (GPT-4) để sinh câu hỏi-đáp án mới, và (3) fine-tune Qwen2.5-4B với QLoRA hoặc Full
Parameter Fine-Tuning (FPFT) tùy tác vụ. Điểm nổi bật là dữ liệu tổng hợp (synthetic data)
được thiết kế theo phương pháp Chain-of-Thought, giúp mô hình học được cách lập luận từng bước.

Chi phí của phương pháp này khá lớn: cần GPU A100 (80~GB) cho pre-training và FPFT, đồng thời
tốn chi phí API để sinh dữ liệu tổng hợp.

\subsection{MinLegal: Phương pháp hai giai đoạn}

Nhóm MinLegal~\cite{minlegal2025} sử dụng chiến lược hai giai đoạn: (1) train LoRA riêng cho
từng tác vụ trên Qwen2.5-4B, (2) merge các LoRA adapter lại rồi tiến hành full fine-tuning.
Phương pháp này cải thiện đáng kể so với giai đoạn 1 nhờ giai đoạn full fine-tuning, nhưng
đòi hỏi tài nguyên tính toán lớn và quy trình quản lý nhiều checkpoint phức tạp.

\subsection{So sánh với phương pháp đề xuất}

Bảng~\ref{tab:compare_methods} tóm tắt sự khác biệt giữa ba phương pháp.

\begin{table}[H]
\centering
\caption{So sánh các phương pháp fine-tuning LegalSLM}
\label{tab:compare_methods}
\begin{tabular}{lccc}
\toprule
\textbf{Tiêu chí} & \textbf{Bosch@AI} & \textbf{MinLegal} & \textbf{Đề xuất} \\
\midrule
Kích thước mô hình & 4B & 4B & \textbf{1.5B} \\
Phần cứng yêu cầu & A100 (80GB) & A100 (80GB) & \textbf{T4 (15GB)} \\
Số giai đoạn train & 3 & 2 & \textbf{1} \\
Dữ liệu tổng hợp & Có (GPT-4) & Không & \textbf{Không} \\
Domain pre-training & Có & Không & \textbf{Không} \\
Thời gian train & Nhiều giờ & Nhiều giờ & \textbf{19,3 phút} \\
\bottomrule
\end{tabular}
\end{table}

% ============================================================
\chapter{Dữ liệu}
\label{chap:data}

\section{Tổng quan Bộ Dữ liệu}

Nghiên cứu sử dụng bộ dữ liệu \textbf{VLSP 2025 LegalSLM Public-Test}, công bố trên HuggingFace
tại địa chỉ \texttt{VLSP2025-LegalSML/Public-Test}. Bộ dữ liệu bao gồm ba tác vụ, mỗi tác vụ
50 mẫu, được thu thập từ các câu hỏi pháp lý thực tế và chú thích bởi nhóm chuyên gia tại
Việt Nam. Bảng~\ref{tab:dataset_overview} tóm tắt thông tin cơ bản.

\begin{table}[H]
\centering
\caption{Thống kê tổng quan bộ dữ liệu}
\label{tab:dataset_overview}
\begin{tabular}{lcccc}
\toprule
\textbf{Tác vụ} & \textbf{Số mẫu} & \textbf{Đầu vào} & \textbf{Đầu ra} & \textbf{Metric} \\
\midrule
NLI & 50 & Văn bản + câu hỏi & Có / Không & Accuracy \\
MCQ & 50 & Câu hỏi + 4 lựa chọn & A / B / C / D & Accuracy \\
Syllogism & 50 & Tình huống pháp lý & Văn bản tự do & ROUGE-L \\
\midrule
\textbf{Tổng} & \textbf{150} & -- & -- & -- \\
\bottomrule
\end{tabular}
\end{table}

\section{Mô tả Chi tiết Từng Tác vụ}

\subsection{Tác vụ NLI --- Suy diễn Ngôn ngữ Tự nhiên}

Tác vụ NLI trong bộ dữ liệu này có tên đầy đủ là \textit{Natural Language Inference for Legal
Citation Usefulness}, phân loại xem một văn bản pháp lý (legal document) có thể dùng để trả
lời một câu hỏi cụ thể hay không.

Mỗi mẫu gồm bốn trường:
\begin{itemize}
  \item \texttt{legal\_document}: Một đoạn trích từ văn bản quy phạm pháp luật, điều khoản
    hoặc nghị định.
  \item \texttt{specific\_question}: Câu hỏi pháp lý mà người dùng cần giải đáp.
  \item \texttt{question}: Câu hỏi meta, thường là: ``Điều luật được cung cấp có thể dùng để
    trả lời câu hỏi trên hay không?''
  \item \texttt{choices} / \texttt{answer}: Nhãn nhị phân $\{0: \text{Có}, 1: \text{Không}\}$.
\end{itemize}

Điều quan trọng cần hiểu rõ: nhãn ``Có'' không có nghĩa là câu trả lời cho
\texttt{specific\_question} là ``Có'', mà có nghĩa là \textit{văn bản pháp lý đủ thông tin
để suy ra câu trả lời}. Đây là điểm dễ gây nhầm lẫn khi thiết kế prompt.

\begin{figure}[H]
\centering
\fbox{\parbox{0.88\textwidth}{
\small
\textbf{legal\_document:} ``Theo Kết luận 83-KL/TW năm 2024, sau năm 2026 khi Bộ Chính trị
ban hành và triển khai hệ thống Danh mục vị trí việc làm, Trung ương sẽ xem xét cải cách
tiền lương đối với 5 bảng lương mới...''\\[4pt]
\textbf{specific\_question:} ``Khi nào sẽ có 5 bảng lương mới theo cải cách tiền lương?''\\[4pt]
\textbf{question:} ``Điều luật được cung cấp có thể dùng để trả lời câu hỏi trên hay không?''\\[4pt]
\textbf{answer:} Có \quad (vì văn bản đề cập đến lộ trình sau năm 2026)
}}
\caption{Ví dụ mẫu tác vụ NLI}
\label{fig:nli_example}
\end{figure}

Phân tích 50 mẫu NLI cho thấy: 30 mẫu có nhãn ``Có'' (60\%) và 20 mẫu có nhãn ``Không'' (40\%).
Sự mất cân bằng nhẹ này cần lưu ý khi đánh giá: một mô hình chỉ đoán ``Có'' toàn bộ đã đạt
60\% accuracy mà không học được gì.

\subsection{Tác vụ MCQ --- Trả lời Câu hỏi Trắc nghiệm}

Mỗi mẫu MCQ gồm một câu hỏi pháp lý và bốn lựa chọn, chỉ một đúng. Câu hỏi bao quát nhiều
lĩnh vực: thuế (Luật Quản lý Thuế, Nghị định 168), lao động (Bộ luật Lao động 2019, Luật BHXH
2024), đất đai (Luật Đất đai 2024), giao thông (Nghị định 168/2024), doanh nghiệp (Luật Doanh
nghiệp 2020), và chứng khoán.

Mức độ khó khá đồng đều --- từ câu hỏi về định nghĩa khái niệm (``Hoàn trả thư tín dụng là
gì?'') đến câu hỏi so sánh, lý giải quy định (``Tại sao cần xác định loại đất dựa trên giấy
tờ pháp lý?''). Câu hỏi có tính nhiễu cao: ba phương án sai thường là các biến thể rất gần
đúng, chỉ khác nhau ở một vài chi tiết số liệu hoặc điều kiện.

\subsection{Tác vụ Syllogism --- Suy luận Tam đoạn luận Pháp lý}

Đây là tác vụ khó nhất và cũng thú vị nhất. Mỗi mẫu là một tình huống pháp lý phức tạp, yêu
cầu mô hình đưa ra phân tích đầy đủ theo cấu trúc tam đoạn luận:
\begin{itemize}
  \item \textbf{Tiền đề lớn (Major premise)}: Điều khoản pháp luật áp dụng
  \item \textbf{Tiền đề nhỏ (Minor premise)}: Tình huống thực tế cụ thể
  \item \textbf{Kết luận (Conclusion)}: Hậu quả pháp lý theo logic suy diễn
\end{itemize}

Câu trả lời mẫu thường dài từ 100 đến 300 từ, có cấu trúc rõ ràng và trích dẫn chính xác điều
khoản pháp luật. Đây là thách thức lớn với mô hình 1.5B vì đòi hỏi đồng thời: (1) ghi nhớ
kiến thức pháp lý, (2) lập luận theo logic, và (3) tạo văn bản có cấu trúc chặt chẽ.

\section{Phân chia Dữ liệu}

Do số lượng mẫu ít (50 mỗi tác vụ), nhóm sử dụng phân chia train/test 80/20 với random seed
cố định (42) để đảm bảo tái tạo. Bảng~\ref{tab:split} tóm tắt.

\begin{table}[H]
\centering
\caption{Phân chia dữ liệu huấn luyện và đánh giá}
\label{tab:split}
\begin{tabular}{lcccc}
\toprule
\textbf{Tác vụ} & \textbf{Train} & \textbf{Test} & \textbf{Tổng} & \textbf{Tỷ lệ Test} \\
\midrule
NLI & 40 & 10 & 50 & 20\% \\
MCQ & 40 & 10 & 50 & 20\% \\
Syllogism & 40 & 10 & 50 & 20\% \\
\midrule
\textbf{Tổng} & \textbf{120} & \textbf{30} & \textbf{150} & 20\% \\
\bottomrule
\end{tabular}
\end{table}

Tập test 10 mẫu mỗi tác vụ là khá nhỏ --- mỗi mẫu đúng/sai đóng góp 10\% vào điểm số.
Điều này đồng nghĩa kết quả có thể biến động lớn theo cách phân chia ngẫu nhiên. Nhóm không
thực hiện cross-validation do hạn chế tài nguyên, nhưng đây là điểm cần lưu ý khi diễn giải
kết quả.

\section{Định dạng Instruction Tuning}

Để huấn luyện theo phong cách instruction tuning, mỗi mẫu được chuyển sang định dạng hội thoại
theo chat template của Qwen2.5 (system + user + assistant). Hình~\ref{fig:prompt_template} minh
họa cấu trúc prompt cho từng tác vụ.

\begin{figure}[H]
\centering
\fbox{\parbox{0.88\textwidth}{
\small\ttfamily
\textbf{[NLI]} \\
\textit{System}: Bạn là một chuyên gia pháp lý Việt Nam... \\
\textit{User}: Văn bản pháp lý: \{legal\_document\}\\
\phantom{x}Câu hỏi cụ thể: \{specific\_question\}\\
\phantom{x}\{question\} \\
\phantom{x}A. Có \quad B. Không\\
\phantom{x}Hãy chọn đáp án đúng (chỉ trả lời bằng chữ cái A hoặc B). \\
\textit{Assistant}: A \hfill $\leftarrow$ chỉ 1 ký tự \\[6pt]
\textbf{[MCQ]} \\
\textit{User}: Câu hỏi pháp luật: \{question\}\\
\phantom{x}A. \{choices[0]\} \quad B. \{choices[1]\}\\
\phantom{x}C. \{choices[2]\} \quad D. \{choices[3]\}\\
\phantom{x}Hãy chọn đáp án đúng (chỉ trả lời bằng chữ cái A, B, C hoặc D). \\
\textit{Assistant}: C \hfill $\leftarrow$ chỉ 1 ký tự \\[6pt]
\textbf{[Syllogism]} \\
\textit{User}: Tình huống pháp lý: \{question\}\\
\phantom{x}Hãy phân tích theo: 1. Điều khoản pháp luật áp dụng (Tiền đề lớn)\\
\phantom{x}2. Áp dụng vào tình huống cụ thể (Tiền đề nhỏ) 3. Kết luận \\
\textit{Assistant}: \{answer\} \hfill $\leftarrow$ văn bản dài
}}
\caption{Template prompt cho từng tác vụ}
\label{fig:prompt_template}
\end{figure}

Chú ý rằng với NLI và MCQ, câu trả lời mẫu (assistant) chỉ là một ký tự (A, B, C, D). Điều
này giúp mô hình học được hành vi ``trả lời ngắn gọn'' và làm đơn giản hóa quá trình parsing
kết quả khi đánh giá.

% ============================================================
\chapter{Phương pháp Đề xuất}
\label{chap:method}

\section{Tổng quan Pipeline}

Hình~\ref{fig:pipeline} mô tả quy trình tổng thể của phương pháp đề xuất. Toàn bộ pipeline
gồm bốn giai đoạn tuần tự: tiền xử lý dữ liệu, thiết lập mô hình với QLoRA, huấn luyện đa
tác vụ, và đánh giá.

\begin{figure}[H]
\centering
\fbox{\parbox{0.85\textwidth}{
\centering
\vspace{0.3cm}
\textbf{Pipeline tổng thể --- QLoRA Multi-task}\\[0.5cm]
\begin{tabular}{c}
$\boxed{\text{Dữ liệu VLSP 2025: NLI (50) + MCQ (50) + Syllogism (50)}}$ \\[0.2cm]
$\downarrow$ \textit{Phân chia train/test 80/20, định dạng instruction tuning} \\[0.2cm]
$\boxed{\text{Tập huấn luyện đa tác vụ (120 mẫu, trộn ngẫu nhiên)}}$ \\[0.2cm]
$\downarrow$ \textit{Load Qwen2.5-1.5B (4-bit NF4) + cấu hình LoRA}$\uparrow$ \\[0.2cm]
$\boxed{\text{QLoRA Fine-tuning: 3 epoch, batch=2, grad\_accum=8}}$ \\[0.2cm]
$\downarrow$ \textit{Lưu LoRA adapter checkpoint} \\[0.2cm]
$\boxed{\text{Đánh giá: Accuracy (NLI/MCQ) + ROUGE-L (Syllogism)}}$
\end{tabular}
\vspace{0.3cm}
}}
\caption{Pipeline tổng thể của phương pháp đề xuất}
\label{fig:pipeline}
\end{figure}

\section{Mô hình Nền: Qwen2.5-1.5B-Instruct}

Qwen2.5-1.5B-Instruct~\cite{qwen25} là phiên bản nhỏ nhất trong dòng mô hình Qwen2.5 do
Alibaba Cloud phát triển và công bố vào năm 2024. Các đặc điểm kỹ thuật chính:

\begin{itemize}
  \item \textbf{Kiến trúc}: Transformer decoder-only, 28 lớp, hidden dim 1536, 12 attention heads
  \item \textbf{Số tham số}: 1,54 tỷ (1.54B)
  \item \textbf{Cửa sổ ngữ cảnh}: 128K token (một trong những context window lớn nhất với mô hình
    ở kích thước này)
  \item \textbf{Đa ngôn ngữ}: Được pre-train trên corpus 7+ nghìn tỷ token bao gồm tiếng Việt,
    tiếng Trung, tiếng Anh và nhiều ngôn ngữ khác
  \item \textbf{Phiên bản}: Dùng phiên bản ``Instruct'' --- đã qua instruction tuning và RLHF
    --- thay vì phiên bản ``Base'' thuần túy, do phù hợp hơn với tác vụ hội thoại/hỏi đáp
\end{itemize}

Lý do chọn Qwen2.5-1.5B thay vì Qwen2.5-4B (phổ biến hơn trong các bài tham chiếu):
\begin{itemize}
  \item Chạy được trên GPU T4 15GB của Colab free tier với lượng tử hóa 4-bit, trong khi 4B
    model cần khoảng 10GB chỉ để load và cần GPU A100 hoặc trả phí Colab Pro cho training
  \item Thời gian fine-tune ngắn hơn ($\approx 20$ phút vs $\approx 2$ giờ)
  \item Vẫn đủ năng lực cho tác vụ phân loại (NLI, MCQ) sau fine-tune
\end{itemize}

\section{Cấu hình QLoRA}

Bảng~\ref{tab:qlora_config} trình bày toàn bộ siêu tham số sử dụng trong thực nghiệm.

\begin{table}[H]
\centering
\caption{Cấu hình QLoRA và siêu tham số huấn luyện}
\label{tab:qlora_config}
\begin{tabular}{lll}
\toprule
\textbf{Nhóm} & \textbf{Tham số} & \textbf{Giá trị} \\
\midrule
\multirow{5}{*}{Quantization} & Bit-width & 4-bit NF4 \\
 & Compute dtype & float16 \\
 & Double quantization & Bật \\
 & -- & -- \\
 & -- & -- \\
\midrule
\multirow{5}{*}{LoRA} & Rank $r$ & 16 \\
 & Alpha $\alpha$ & 32 ($= 2r$) \\
 & Dropout & 0.05 \\
 & Bias & none \\
 & Target modules & q, k, v, o, gate, up, down\_proj \\
\midrule
\multirow{8}{*}{Training} & Learning rate & 2e-4 \\
 & Batch size per device & 2 \\
 & Gradient accumulation & 8 \\
 & Effective batch size & 16 \\
 & Epochs & 3 \\
 & Max sequence length & 1024 token \\
 & Optimizer & paged\_adamw\_8bit \\
 & LR scheduler & cosine \\
 & Warmup ratio & 0.05 \\
\bottomrule
\end{tabular}
\end{table}

Một số lưu ý về các lựa chọn siêu tham số:

\textbf{Target modules:} Nhóm áp dụng LoRA cho cả 7 nhóm projection (q, k, v, o, gate, up,
down), không chỉ q và v như trong paper gốc của LoRA. Cách này tăng số tham số trainable nhưng
vẫn giữ tỷ lệ huấn luyện ở mức thấp ($\approx 1.52\%$).

\textbf{fp16=False, bf16=False:} Không dùng Mixed Precision Training của PyTorch. Lý do: khi mô
hình nền đã được lượng tử hóa 4-bit bởi bitsandbytes, việc bật fp16 sẽ kích hoạt GradScaler
của PyTorch AMP, dẫn đến lỗi ``not implemented for BFloat16'' trong quá trình rescale gradient.
Bitsandbytes đã tự xử lý precision nội bộ nên không cần AMP bên ngoài.

\textbf{paged\_adamw\_8bit:} Optimizer 8-bit với paged memory --- hỗ trợ trao đổi optimizer
states ra CPU RAM khi GPU VRAM đầy, đặc biệt hữu ích trên T4 15GB.

\section{Chiến lược Huấn luyện Đa Tác vụ}

Sau khi định dạng dữ liệu thành instruction tuning format, 120 mẫu training (40 NLI + 40 MCQ
+ 40 Syllogism) được trộn ngẫu nhiên thành một tập duy nhất. Quá trình huấn luyện chạy một
lần duy nhất trên tập này qua 3 epoch.

Điều này khác căn bản với chiến lược của MinLegal (train riêng lẻ từng tác vụ). Ưu điểm của
cách tiếp cận của nhóm:

\begin{enumerate}
  \item \textbf{Một LoRA adapter duy nhất} phục vụ được cả ba tác vụ --- không cần chuyển đổi
    model khi thay tác vụ.
  \item \textbf{Tăng hiệu quả data}: Với 120 mẫu trong một lần train thay vì 3 lần train 40
    mẫu, mô hình nhìn thấy toàn bộ dữ liệu trong mỗi epoch.
  \item \textbf{Regularization qua đa dạng}: Sự khác biệt giữa các tác vụ hoạt động như một
    hình thức regularization tự nhiên.
\end{enumerate}

Trong quá trình huấn luyện, loss được tính trên toàn bộ chuỗi token của assistant response.
Với NLI và MCQ chỉ có 1 token response (``A'', ``B''...), loss chủ yếu đến từ Syllogism. Điều
này có thể giải thích tại sao loss cuối khá cao (0.9357) --- nó phản ánh độ khó của tác vụ
sinh văn bản tự do hơn là tác vụ phân loại.

\section{Quy trình Đánh giá}

\textbf{NLI và MCQ:} Với mỗi mẫu test, mô hình được cho prompt và sinh tối đa 16 token. Token
đầu tiên trong output được trích xuất và đối chiếu với đáp án đúng. Nếu token đầu tiên không
phải là chữ cái hợp lệ (A--D), thuật toán tìm chữ cái đầu tiên xuất hiện trong 10 ký tự đầu
của output. Accuracy = số câu đúng / tổng số câu.

\textbf{Syllogism:} Mô hình sinh tối đa 512 token. ROUGE-L F1 được tính giữa output sinh ra
và đáp án tham chiếu, sử dụng thư viện \texttt{rouge-score} của Google, với \texttt{use\_stemmer=False}
(không stemming do đặc thù tiếng Việt).

% ============================================================
\chapter{Thực nghiệm và Kết quả}
\label{chap:results}

\section{Môi trường Thực nghiệm}

\begin{table}[H]
\centering
\caption{Cấu hình môi trường thực nghiệm}
\label{tab:env}
\begin{tabular}{ll}
\toprule
\textbf{Thành phần} & \textbf{Chi tiết} \\
\midrule
Nền tảng & Google Colab (free tier) \\
GPU & NVIDIA Tesla T4, 15 GB VRAM \\
CPU & Intel(R) Xeon(R) (2 vCPU) \\
RAM & 12.7 GB \\
Hệ điều hành & Ubuntu 22.04 \\
Python & 3.10 \\
PyTorch & 2.4.1+cu121 \\
Transformers & 4.46.x \\
PEFT & 0.13.x \\
TRL & 0.12.x \\
BitsAndBytes & 0.45.x \\
\midrule
Mô hình nền & Qwen/Qwen2.5-1.5B-Instruct \\
Số tham số tổng & $\approx 1{,}54\text{B}$ \\
Số tham số trainable & $\approx 23{,}4\text{M}$ (1,52\%) \\
\bottomrule
\end{tabular}
\end{table}

\section{Kết quả Chính}

Bảng~\ref{tab:main_results} trình bày kết quả so sánh giữa baseline zero-shot và phương pháp
QLoRA đa tác vụ trên tập test (30 mẫu: 10 mỗi tác vụ).

\begin{table}[H]
\centering
\caption{Kết quả đánh giá trên tập test (30 mẫu)}
\label{tab:main_results}
\begin{tabular}{lccc}
\toprule
\multirow{2}{*}{\textbf{Phương pháp}} & \textbf{NLI} & \textbf{MCQ} & \textbf{Syllogism} \\
\cmidrule(lr){2-2}\cmidrule(lr){3-3}\cmidrule(lr){4-4}
 & Accuracy & Accuracy & ROUGE-L \\
\midrule
Zero-shot (Qwen2.5-1.5B) & 0,6667 & 0,7333 & 0,4230 \\
\textbf{QLoRA Multi-task (đề xuất)} & \textbf{0,9333} & \textbf{0,8667} & \textbf{0,5066} \\
\midrule
Cải thiện tuyệt đối & +0,2666 & +0,1334 & +0,0836 \\
Cải thiện tương đối & \textbf{+40,0\%} & \textbf{+18,2\%} & \textbf{+19,8\%} \\
\bottomrule
\end{tabular}
\end{table}

\textit{Chú thích: BERTScore không được tính cho Syllogism do hạn chế bộ nhớ GPU. Kết quả
được đo trên tập test cố định (random\_seed=42).}

\section{Phân tích Quá trình Huấn luyện}

\begin{table}[H]
\centering
\caption{Thống kê quá trình huấn luyện}
\label{tab:training_stats}
\begin{tabular}{ll}
\toprule
\textbf{Chỉ số} & \textbf{Giá trị} \\
\midrule
Training loss cuối (epoch 3) & 0,9357 \\
Thời gian huấn luyện & 19,3 phút \\
Số tham số được huấn luyện & $\approx 23{,}4$M (1,52\% tổng) \\
Tổng số tham số & $\approx 1{,}54$B \\
Số mẫu train & 120 (40 mỗi tác vụ) \\
Effective batch size & 16 (batch=2, grad\_accum=8) \\
Tổng số bước (steps) & $120 / 16 \times 3 = 22{,}5 \approx 23$ steps \\
\bottomrule
\end{tabular}
\end{table}

Training loss cuối là 0,9357 --- con số khá cao nếu so với các task đơn lẻ thông thường (thường
$<0.5$). Điều này phản ánh độ khó tổng hợp của tập dữ liệu đa tác vụ, đặc biệt do Syllogism
yêu cầu sinh văn bản dài và phức tạp. Khi chỉ nhìn vào NLI và MCQ (single token response),
model thực tế đã học rất tốt (NLI 93,33\%, MCQ 86,67\%), cho thấy loss cao chủ yếu đến từ
Syllogism.

\section{Phân tích Định tính}

\subsection{Ví dụ 1: Tác vụ NLI --- Mô hình được cải thiện}

\begin{figure}[H]
\centering
\fbox{\parbox{0.88\textwidth}{
\small
\textbf{Văn bản pháp lý}: ``Theo Kết luận 83-KL/TW năm 2024, sau năm 2026 khi Bộ Chính trị
ban hành và triển khai hệ thống Danh mục vị trí việc làm, Trung ương sẽ xem xét cải cách
tiền lương đối với 5 bảng lương mới...''\\[4pt]
\textbf{Câu hỏi cụ thể}: ``Khi nào sẽ có 5 bảng lương mới theo cải cách tiền lương?''\\[4pt]
\textbf{Câu hỏi meta}: ``Điều luật có thể dùng để trả lời câu hỏi trên hay không?''\\[4pt]
\textbf{Zero-shot}: ``B'' (Không) --- \textit{\color{red} sai}\\
\textbf{QLoRA}: ``A'' (Có) --- \textit{\color{green!50!black} đúng} \\[4pt]
\textit{Phân tích}: Zero-shot hiểu nhầm câu hỏi --- mô hình nghĩ câu hỏi hỏi về thời điểm
cụ thể (ngày tháng năm) nên thấy văn bản chỉ nói ``sau năm 2026'' và cho rằng không đủ thông
tin. Sau fine-tune, mô hình hiểu rằng ``sau năm 2026'' là câu trả lời hợp lệ cho ``Khi nào''.
}}
\caption{Ví dụ NLI: so sánh zero-shot và sau fine-tune}
\label{fig:nli_qual}
\end{figure}

\subsection{Ví dụ 2: Tác vụ MCQ --- Cả hai đều đúng}

\begin{figure}[H]
\centering
\fbox{\parbox{0.88\textwidth}{
\small
\textbf{Câu hỏi}: ``Theo quy định pháp luật, người nộp thuế có nghĩa vụ gì liên quan đến
ghi mã số thuế trên hóa đơn?''\\[4pt]
\textbf{Lựa chọn}: A. Phải ghi trong mọi trường hợp / B. Chỉ ghi khi bên mua yêu cầu /
\textbf{C. Ghi mã số thuế được cấp trên hóa đơn, trừ bên mua không có} / D. Không bắt buộc\\[4pt]
\textbf{Zero-shot}: ``C'' --- \textit{\color{green!50!black} đúng}\\
\textbf{QLoRA}: ``C'' --- \textit{\color{green!50!black} đúng}\\[4pt]
\textit{Phân tích}: Câu hỏi này không phân biệt được hai phương pháp. Kết quả tốt trên MCQ
(73\% zero-shot) cho thấy Qwen2.5-1.5B đã có kiến thức pháp lý nền tảng; fine-tune giúp cải
thiện thêm những câu khó hơn.
}}
\caption{Ví dụ MCQ: cả zero-shot và QLoRA đều đúng}
\label{fig:mcq_qual}
\end{figure}

\subsection{Ví dụ 3: Tác vụ MCQ --- Chỉ QLoRA đúng}

Một ví dụ điển hình cho sự cải thiện của fine-tune là các câu hỏi liên quan đến số liệu cụ
thể trong văn bản quy phạm:

\begin{quote}
\textit{Câu hỏi}: ``Mức phạt tiền khi lái ô tô vượt đèn đỏ gây tai nạn theo Nghị định
168/2024 là bao nhiêu?''\\
\textit{Zero-shot}: Chọn B (50--60 triệu đồng) --- \textbf{sai}\\
\textit{QLoRA}: Chọn A (50--70 triệu đồng) --- \textbf{đúng}
\end{quote}

Câu hỏi về con số cụ thể trong văn bản pháp luật ban hành năm 2024 nằm ngoài kiến thức
pre-training của mô hình (cutoff~2023). Fine-tune trên dữ liệu gần đây hơn giúp mô hình
nhớ các giá trị số này.

\section{Phân tích Lỗi}

\subsection{Lỗi của mô hình sau fine-tune}

Phân tích các mẫu mà QLoRA vẫn trả lời sai cho thấy một số mẫu lỗi hệ thống:

\textbf{NLI (1/10 câu sai)}: Mô hình vẫn còn nhầm lẫn ở một số trường hợp văn bản pháp lý
dài (>500 từ) chứa nhiều điều khoản liên quan nhưng không trực tiếp trả lời câu hỏi cụ thể.
Mô hình có xu hướng đánh giá ``Có'' khi thấy từ khóa trùng với câu hỏi, ngay cả khi văn bản
không thực sự cung cấp câu trả lời.

\textbf{MCQ (2/10 câu sai)}: Các lỗi thường gặp ở câu hỏi ``So sánh'' (so sánh hai quy định
khác nhau) và câu hỏi về quy trình nhiều bước. Mô hình 1.5B có giới hạn trong việc theo dõi
nhiều ràng buộc logic đồng thời.

\textbf{Syllogism}: ROUGE-L 0.5066 ngụ ý mô hình tạo ra khoảng 50\% nội dung trùng lặp với
đáp án chuẩn. Quan sát định tính cho thấy mô hình thường đúng về cấu trúc (tiền đề lớn/nhỏ/kết
luận) nhưng thiếu chính xác về trích dẫn điều khoản cụ thể.

\subsection{So sánh với giới hạn lý thuyết}

Với tập test chỉ 10 mẫu mỗi tác vụ, độ chính xác tối đa có thể là 100\% (10/10). Phương pháp
của nhóm đạt 9/10 cho NLI và 8/10 cho MCQ. Để đạt 10/10, cần dữ liệu huấn luyện đầy đủ hơn
hoặc mô hình lớn hơn.

\section{Thảo luận}

\subsection{Điểm mạnh của phương pháp}

Kết quả thực nghiệm khẳng định rằng QLoRA đa tác vụ là hướng đi hiệu quả và phù hợp với
điều kiện tài nguyên hạn chế:

\begin{itemize}
  \item \textbf{NLI +40\%}: Cải thiện lớn nhất thuộc về tác vụ phân loại nhị phân. Điều này
    có thể giải thích bởi bản chất của tác vụ: mô hình chỉ cần học phân biệt hai loại tình
    huống (văn bản có/không liên quan đến câu hỏi), và instruction tuning đã dạy mô hình tập
    trung vào đúng khía cạnh cần xem xét.
  \item \textbf{Chi phí cực thấp}: 19,3 phút và 0 đô la (Colab free tier). Đây là điểm quan
    trọng nhất từ góc độ thực tế: sinh viên hoặc doanh nghiệp nhỏ có thể thực hiện loại nghiên
    cứu này mà không cần đầu tư phần cứng.
  \item \textbf{Đơn giản}: Không cần pipeline phức tạp, không cần dữ liệu tổng hợp, không cần
    domain pre-training. Chỉ cần khoảng 100 mẫu mỗi tác vụ.
\end{itemize}

\subsection{Hạn chế}

\begin{itemize}
  \item \textbf{Tập test quá nhỏ}: 10 mẫu/tác vụ khiến kết quả dễ bị ảnh hưởng bởi cách phân
    chia ngẫu nhiên. Một phân chia khác có thể cho kết quả dao động $\pm 10\%$.
  \item \textbf{Mô hình nhỏ}: So với Bosch@AI và MinLegal dùng 4B, kích thước 1.5B giới hạn
    khả năng lý luận phức tạp, đặc biệt cho Syllogism. Không có con số so sánh trực tiếp
    trên cùng tập test, nhưng khả năng cao 4B sẽ vượt trội hơn đáng kể ở tác vụ này.
  \item \textbf{ROUGE-L không đủ cho Syllogism}: Điểm 0.5066 nghe có vẻ thấp nhưng chưa phản
    ánh đầy đủ chất lượng thực tế. Cần LLM-as-judge (dùng GPT-4 đánh giá) để có nhận xét
    chính xác hơn về chất lượng lập luận pháp lý.
  \item \textbf{Không có few-shot baseline}: Nhóm dự định đánh giá few-shot (3-shot) nhưng
    không thực hiện được do hạn chế thời gian.
\end{itemize}

% ============================================================
\chapter*{Kết luận và Hướng phát triển}
\label{chap:conclusion}
\addcontentsline{toc}{chapter}{\nameref{chap:conclusion}}

\section*{Tóm tắt kết quả}

Nghiên cứu này đã trình bày và kiểm chứng thực nghiệm một phương pháp tinh chỉnh mô hình
ngôn ngữ nhỏ (Qwen2.5-1.5B) cho ba tác vụ suy luận pháp lý tiếng Việt bằng kỹ thuật QLoRA
và chiến lược huấn luyện đa tác vụ. Kết quả đạt được:

\begin{itemize}
  \item NLI: 66,67\% $\to$ \textbf{93,33\%} (+40,0\%)
  \item MCQ: 73,33\% $\to$ \textbf{86,67\%} (+18,2\%)
  \item Syllogism ROUGE-L: 0,4230 $\to$ \textbf{0,5066} (+19,8\%)
  \item Thời gian huấn luyện: \textbf{19,3 phút} trên GPU T4 miễn phí
\end{itemize}

Những con số này cho thấy rằng ngay cả với mô hình nhỏ, dữ liệu ít và tài nguyên hạn chế,
instruction tuning có thể tạo ra cải thiện rõ rệt và đo lường được.

\section*{Hướng phát triển}

Có một số hướng có thể mở rộng từ kết quả này:

\textbf{(1) Tăng quy mô dữ liệu:} Hiện tại chỉ dùng 40 mẫu/tác vụ. Tăng lên 200--500 mẫu
qua data augmentation (dùng LLM lớn hơn để paraphrase các câu hỏi có sẵn) nhiều khả năng
sẽ cải thiện đáng kể, đặc biệt cho Syllogism.

\textbf{(2) Thử mô hình lớn hơn:} Qwen2.5-3B hoặc Qwen2.5-4B vẫn có thể train được trên T4
với QLoRA. Với 4B, khả năng lý luận tốt hơn sẽ giúp cải thiện Syllogism đáng kể.

\textbf{(3) Curriculum learning:} Thay vì trộn ngẫu nhiên, train theo thứ tự từ tác vụ đơn
giản đến phức tạp (NLI $\to$ MCQ $\to$ Syllogism) có thể giúp mô hình học tuần tự và
tránh interference.

\textbf{(4) Đánh giá LLM-as-judge:} Với Syllogism, dùng GPT-4 hoặc Gemini Pro để đánh giá
chất lượng lập luận thay cho ROUGE-L sẽ cho kết quả đáng tin cậy và có ý nghĩa hơn nhiều.

\textbf{(5) Domain-specific vocabulary:} Thêm các thuật ngữ pháp lý tiếng Việt vào tokenizer
(vocabulary expansion) có thể giúp mô hình xử lý hiệu quả hơn các từ đặc thù như tên điều
khoản, số nghị định, v.v.

% ============================================================
\bibliographystyle{plain}
\begin{thebibliography}{99}

\bibitem{bosch2025}
Tran Minh Quang, Nguyen Xuan Phi, Nguyen Van Tai, Phan Minh Toan, Dang Van Thin.
\textit{Bosch@AI\_Team at LegalSML 2025: Vietnamese Legal Small Language with Domain Adaptation
and Aspect-based Data Synthesis}.
Proceedings of VLSP Workshop 2025.

\bibitem{minlegal2025}
Ngo Dinh Luan, Kiet Nguyen Van.
\textit{MinLegal at VLSP2025-LegalSLM: A Two-Stage LoRA-to-Full Fine-tuning Approach for
Vietnamese Legal Small Language Models}.
Proceedings of VLSP Workshop 2025.

\bibitem{vlsp2025}
VLSP 2025 Organizing Committee.
\textit{VLSP 2025 Shared Task on Vietnamese Legal Small Language Models (LegalSLM)}.
\url{https://vlsp.org.vn/vlsp2025/eval/legalslm}, 2025.

\bibitem{vlqa2025}
Nguyen et al.
\textit{VLQA: A Comprehensive Vietnamese Legal Question Answering Dataset}.
2025.

\bibitem{qwen25}
Qwen Team.
\textit{Qwen2.5 Technical Report}.
Alibaba Cloud, 2024.
\url{https://arxiv.org/abs/2412.15115}

\bibitem{lora}
Edward J. Hu, Yelong Shen, Phillip Wallis, Zeyuan Allen-Zhu, Yuanzhi Li, Shean Wang,
Lu Wang, Weizhu Chen.
\textit{LoRA: Low-Rank Adaptation of Large Language Models}.
International Conference on Learning Representations (ICLR), 2022.

\bibitem{qlora}
Tim Dettmers, Artidoro Pagnoni, Ari Holtzman, Luke Zettlemoyer.
\textit{QLoRA: Efficient Finetuning of Quantized LLMs}.
Neural Information Processing Systems (NeurIPS), 2023.

\bibitem{ouyang2022}
Long Ouyang, Jeffrey Wu, Xu Jiang, et al.
\textit{Training language models to follow instructions with human feedback (InstructGPT)}.
NeurIPS, 2022.

\bibitem{mtl_survey}
Yu Zhang, Qiang Yang.
\textit{A Survey on Multi-Task Learning}.
IEEE Transactions on Knowledge and Data Engineering, 34(12), 2022.

\bibitem{flan2022}
Hyung Won Chung, Le Hou, Shayne Longpre, et al.
\textit{Scaling Instruction-Finetuned Language Models (FLAN-T5)}.
arXiv:2210.11416, 2022.

\end{thebibliography}

\end{document}
